Suppose $\ell\mid k_n$. \cite{KY} Prop.~4.1 gives a saturated component $f_0$ (Lemma~E). The accepted line on $f_0$ names a vector $a\in L_{f_0}\setminus\ell R_n$ with $v=\ell^{-1}H_na\ne0$; by Lemma~E, $v$ is the log vector of $r_a\in\RE$, $r_a\ne\pm1$. A RHO line asserts $|v|<L_n$, contradicting Lemma~A with \cite{MO16} Lemma~2.5(1); a T line asserts $\sum_j\sigma^j(r_a)^2=\sum_j e^{2y_j}<\bar T_n$, where $\bar T_n=33\cdot2^n$ for $n\ge3$ and $\bar T_n=17\cdot2^n$ for $n=2$ (the verifier recomputes $\bar T_n$ from $n$), contradicting the trace floor $\operatorname{Tr}(r_a^2)\ge\bar T_n$ of \cite{KY} Thm.~2.3 for $n\ge3$ and of \cite{MO3} Prop.~6.6 for $n=2$ (the uniform constant $33\cdot2^n$ would not give a contradiction at $n=2$ for $68\le T<132$). Hence $\ell\nmid k_n$.
